% ACM CI 2026 Poster submission (non-archival, non-anonymous)
% Main body must be <= 2 pages; references may extend beyond.
\documentclass[manuscript,screen]{acmart}

\AtBeginDocument{%
  \providecommand\BibTeX{{Bib\TeX}}}

\setcopyright{none}
\copyrightyear{2026}
\acmYear{2026}
\acmDOI{}
\acmISBN{}
\acmConference[CI '26]{ACM Collective Intelligence 2026}{September 27--30, 2026}{Virginia Tech, USA}
\acmBooktitle{Posters of the 14th ACM Collective Intelligence Conference (CI '26), September 27--30, 2026, Virginia Tech, USA}

\ccsdesc[500]{Human-centered computing~Collaborative and social computing systems and tools}
\ccsdesc[300]{Computing methodologies~Multi-agent systems}
\ccsdesc[300]{Applied computing~Economics}

\title{Agents Unite: Crowdsourced LLM Agents as a Git-Native Financial Memory}

\author{Parag Paul}
\affiliation{%
  \institution{RIG AI}
  \city{Remote}
  \country{USA}
}
\email{parag@rigai.co}

\renewcommand{\shortauthors}{Paul}

\begin{abstract}
Large language model (LLM) agents can synthesize market narratives at scale, but solo deployments waste tokens, discard history, and miss timezone-specific information. We present \emph{Agents Unite}, an architecture that treats GitHub as a collective memory substrate for distributed financial intelligence. Each participant runs a scheduled local agent that researches one deterministically assigned ticker per day, submits a structured report as a pull request, and optionally serves as a verifier or consensus writer. Raft-style leader election resolves hourly write conflicts; reputation-weighted aggregation merges independent reports into a dated ledger. The design applies collective intelligence ideas (task decomposition, diversity, and consensus) without blockchain infrastructure: Git commits, CI validation, and CODEOWNERS provide auditability. Per-operator inference cost drops from hundreds of dollars per day for solo full-universe coverage to roughly five cents per contributor per day; readers fork the archive at zero marginal inference cost. This poster presents an early proof-of-concept architecture and invites feedback on assignment, verification, and herding mitigation before a public benchmark evaluation.
\end{abstract}

\keywords{collective intelligence, human-AI collaboration, multi-agent systems, crowdsourcing, financial data}

\begin{document}
\maketitle

\section{Motivation and Collective Intelligence Framing}
Financial markets generate narrative, sentiment, and price information faster than any single analyst or agent can absorb. Recent LLM agents help, but three failures persist in solo deployments: \emph{ephemeral memory} (yesterday's analysis must be regenerated at full cost), \emph{budget exhaustion} (covering thousands of tickers daily costs hundreds to thousands of dollars in API tokens), and \emph{temporal blind spots} (a single timezone misses overnight filings and foreign-market chatter).

Agents Unite reframes market research as a \textbf{crowd intelligence} problem rather than a single-agent scaling problem~\cite{agentsunite}. Instead of one operator researching the entire universe, many contributors each research one ticker per cycle. Assignment is deterministic (contributors cannot cherry-pick popular symbols), and roles (research, verify, consensus) are assigned at run time by a scheduled local runner. The system is designed for human-AI complementarity: humans may contribute manually or run agents with different model providers (hosted APIs, local models, or coding-agent adapters), while CI and peer verifiers provide lightweight quality control.

This work is affiliated with the \textbf{collective intelligence (CI)} track: it studies how decentralized agents, diverse information sources, and consensus mechanisms combine into a longitudinal memory that no individual node could economically maintain alone.

\section{System Design}
Figure~\ref{fig:pipeline} summarizes the pipeline. There is no central application server; GitHub is the database.

\begin{figure}[h]
  \centering
  \fbox{\parbox{0.92\linewidth}{\centering\vspace{0.4em}
    Assign $\rightarrow$ Research $\rightarrow$ Validate $\rightarrow$ PR $\rightarrow$ CI $\rightarrow$ Verify $\rightarrow$ Consensus $\rightarrow$ \texttt{data/DATE/TICKER/}
  \vspace{0.4em}}}
  \caption{End-to-end ingestion: local agents produce reports; GitHub CI and peer verifiers gate merges.}
  \Description{A linear pipeline from assignment through research, validation, pull request, CI, verification, consensus, and canonical data storage.}
  \label{fig:pipeline}
\end{figure}

\textbf{Deterministic assignment.} Ticker $T_k$ is chosen by hashing the date and contributor identity modulo universe size $|T|$. A parallel lottery assigns roles, adjusted by a CI-generated coverage manifest that prioritizes under-covered sectors, geographies, and modalities.

\textbf{Verification without expensive re-research.} Verifiers inspect public pull requests using smaller models and deterministic consistency checks on cited sources, escalating uncertain claims to yes/no or cross-encoder natural language inference checks. Participants do not pre-declare as verifiers; the schedule decides.

\textbf{Hourly leader election.} For Phase~3 hourly shards, multiple pull requests may arrive for the same \texttt{(date, ticker, hour)}. After the hour window closes, CI recomputes
\[
  \text{leader}=\arg\min_i \mathrm{SHA256}(\text{date}:\text{hour}:\text{ticker}:\text{hash}_i)
\]
among schema-valid submissions. Only the leader merges into the canonical path. Provenance is recorded in the pull request, merge commit, and per-ticker metadata.

\textbf{Consensus.} Validated sentiment scores pass through median aggregation with median absolute deviation outlier rejection, then a reputation-weighted median. Source-diversity weighting is proposed to mitigate herding when contributors cite the same wire story~\cite{da2020}.

\section{Economics and Early Evidence}
Table~\ref{tab:econ} contrasts solo regeneration with crowd contribution. The gain is per-operator burden reduction: aggregate work still happens, but it is spread across contributors and reused by readers who fork the ledger.

\begin{table}[h]
  \caption{Illustrative per-operator annual costs (July 2026 API list rates).}
  \label{tab:econ}
  \begin{tabular}{@{}lr@{}}
    \toprule
    \textbf{Path} & \textbf{Cost} \\
    \midrule
    Solo full-universe (budget) & $\sim$\$73k/yr \\
    Solo full-universe (mid) & $\sim$\$219k/yr \\
    \textbf{Agents Unite contributor} & \textbf{$\sim$\$18/yr} \\
    \textbf{Agents Unite reader} & \textbf{\$0 inference} \\
    \bottomrule
  \end{tabular}
\end{table}

External work on federated inference-tree sharing suggests pooled agent memory can sharply outperform single-node caches at fixed false-accept rates~\cite{folklore2026}; Agents Unite aims to test an analogous effect on ticker-day financial queries once a public benchmark exists. Empirical validation (sentiment factor backtests, ledger reuse, verification cost, and adversarial probes) is planned for a follow-on paper after benchmark release.

\section{Poster Contribution and Discussion Questions}
This poster presents an \textbf{early architecture} for turning fragmented agent labor into a compounding public memory. We seek feedback on four open design choices:
\begin{enumerate}
  \item Whether deterministic assignment plus source-diversity weighting is sufficient to limit herding~\cite{da2020}.
  \item How much verification can be delegated to cheap binary checks versus full agentic re-research.
  \item Whether Git-native coordination is adequate at multi-thousand-report-per-day scale without hybrid columnar snapshots.
  \item What governance model, if any, should steward a cross-asset narrative ledger as participation grows.
\end{enumerate}

Outputs are structured observations, not investment advice. A companion architecture paper and conference poster provide additional detail; a live demo repository is under active development.

\begin{acks}
We thank early contributors who tested schema and CI workflows.
\end{acks}

\bibliographystyle{ACM-Reference-Format}
\bibliography{../paper/references}

\end{document}
